% Clean LaTeX contribution for the parent; statement numbers to be assigned
% in the master paper. Complete extended audit is work/ghost_audit.md.
\subsection{Ordinary completed pairs and the exact limit question}
We use the power completion
\[
 \mathcal K=\varprojlim_N\mathbb Q^\times/(\mathbb Q^\times)^N,
\]
with its inverse-limit topology, and write \(\mathcal G_{\rm ord}\)
for the pairs whose localization at each finite prime is
\((\kappa_p(t_p),\kappa_p(1-t_p))\), with
\(t_p\in\mathbb Q_p\setminus\{0,1\}\).
Unique factorization gives the actual topological isomorphism
\[
 \mathcal K\simeq\{\pm1\}\times
 \left\{(a_q)_q\in\prod_q\widehat{\mathbb Z}:
  \text{for each }N,\ a_q\equiv0\pmod N
  \text{ outside a finite set}\right\}.
\]
The topology on the right is the one induced by the simultaneous
finite-support power quotients. At level \(N\), this is just the
factorization of \(\mathbb Q^\times\) into its sign and prime
powers; compatible level classes prove both directions of the
isomorphism. The map to the \(q\)-coordinate is denoted \(v_q\).
For an ordinary pair all these coordinates are ordinary integers,
embedded diagonally in \(\widehat{\mathbb Z}\).

\begin{result}{C.1.1}{Proved}{Finite-support reconstruction; confirms SECTIONS A.32}
\label{prop:ghost-finite-support}
Every pair in \(\mathcal G_{\rm ord}\) for which the valuation
vector of at least one of \(A\), \(B\), or \(A/B\) is
finitely supported has the form
\[
       (A,B)=(\kappa(t),\kappa(1-t)),
       \qquad t\in\mathbb Q\setminus\{0,1\},
\]
for a unique \(t\), and \(t_p=t\) at every finite prime.
\end{result}
\begin{proof}
The local Kummer map \(\mathbb Q_p^\times\to
\varprojlim_N\mathbb Q_p^\times/(\mathbb Q_p^\times)^N\)
is injective. The valuation of an element in its kernel is divisible
by every positive integer. Its valuation is therefore zero; its
finite unit-torsion part is killed by its exponent, and its principal
unit part lies in the intersection of all \(p\)-power multiples of
\(\mathbb Z_p\), which is zero. At two use the subgroup
\(1+4\mathbb Z_2\).

For finite support \(S\) of \(A\), its sign \(\epsilon\) and
integer valuations give the rational number
\(t=\epsilon\prod_{q\in S}q^{v_q(A)}\). The displayed
isomorphism gives \(A=\kappa(t)\). Local injectivity gives
\(t_p=t\), so \(t\ne0,1\). All localizations of
\(B/\kappa(1-t)\) are trivial. Their valuations kill every prime
coordinate. The remaining sign is detected at three, where
\(-1\) is not a square. Hence \(B=\kappa(1-t)\). Interchanging
the two coordinates gives the other case, and local injectivity
gives uniqueness.

For finite support of \(A/B\), the same valuation isomorphism
gives \(A/B=\kappa(u)\) for a rational \(u\ne0\). Local
injectivity gives \(t_p/(1-t_p)=u\). The value \(u=-1\)
would imply \(0=-1\), so \(u\ne-1\), and
\(t_p=u/(1+u)\) at every place. The preceding reconstruction
applied to the resulting rational local classes gives the
asserted pair.
\end{proof}

Fix a rational tangential base point, write
\(1\to\Pi\to E\to G_{\mathbb Q}\to1\), and take the closed
lower central series \(\Gamma_1\Pi=\Pi\),
\(\Gamma_{r+1}\Pi=\overline{[\Pi,\Gamma_r\Pi]}\).
Set \(E_r=E/\Gamma_{r+1}\Pi\). A compatible family of literal
continuous sections of the maps \(E_r\to G_{\mathbb Q}\) is
exactly a continuous section into
\[
 E_{\rm nil}=\varprojlim_rE_r
       =E/\bigcap_r\Gamma_{r+1}\Pi.
\]
The equality follows from compactness of the nested cosets in
\(E\); coordinatewise continuity proves the section comparison.
The geometric quotient is canonically
\(\Pi^{\rm nil}=\prod_\ell\Pi^{(\ell)}\): every finite
\(\ell\)-group is nilpotent, and every finite nilpotent group
is the product of its Sylow groups. These identify the same finite
quotients on both sides and hence the actual profinite groups.

\begin{result}{C.1.2}{Proved}{Infinite ramification at every coefficient prime}
\label{prop:ghost-infinite-ramification}
For an ordinary pair with infinite support, its abelian
\(\ell\)-primary Kummer component ramifies at infinitely many
rational primes, for every \(\ell\). Its abelian component at
no \(\ell\) is the component of a rational point or a cuspidal
section.
For every finite set of rational primes \(S\), the vectors
\((v_q(A),v_q(B))\), \(q\notin S\), span
\(\mathbb Q_\ell^2\) for every \(\ell\). In particular the
abelian geometric inertia images outside any finite set
containing \(\ell\) generate an open submodule of
\(\mathbb Z_\ell(1)^2\).
\end{result}
\begin{proof}
At \(q\ne\ell\), write an ordinary local coordinate as
\(q^a u\). Roots of the residue of \(u\) lie in unramified
extensions, and the principal \(q\)-units have unique
\(\ell^n\)-th roots by Hensel's lemma. Thus the restriction of
its \(\ell\)-primary Kummer cocycle to inertia is
\(a\,t_\ell\), where \(t_\ell:I_q\to\mathbb Z_\ell(1)\)
is the surjective tame character. A nonzero ordinary integer stays
nonzero in \(\mathbb Z_\ell\). Every prime in the infinite
support, apart from \(\ell\) itself, therefore ramifies.

At any ordinary local point the nonzero valuation vector has
one of the three forms \((a,0)\), \((0,b)\), or \((c,c)\),
with \(a,b>0\) and \(c<0\). The three corresponding lines
are pairwise independent over \(\mathbb Q_\ell\). A span of
rank at most one outside \(S\) would force all such vectors
onto one of these lines. Respectively, \(B\), \(A\), or
\(A/B\) would then have finite support, contradicting
Proposition~\ref{prop:ghost-finite-support}. Thus the rank is
two. Two independent ordinary integer vectors already generate
a finite-index \(\mathbb Z_\ell\)-submodule, and the tame
character is surjective, giving the stated assertion about the
actual inertia maps. This asserts openness only in the abelian
geometric quotient.

Equality with the abelian image of a rational point gives
\(v_q(A)=v_q(t)\) in \(\mathbb Z_\ell\). Injectivity of
\(\mathbb Z\to\mathbb Z_\ell\) gives equality as ordinary
integers and hence finite support, already excluded. A cusp at
zero or one makes one valuation vector zero. A cusp at infinity
makes the two ordinary valuation vectors equal. Then \(A/B\)
has only its sign coordinate. Ordinary local injectivity gives
\(t_q/(1-t_q)=1\) or \(-1\); the former gives the rational pair
at \(1/2\), and the latter gives the impossible equality
\(1=0\). This excludes all three cusps.
\end{proof}

\begin{result}{C.2.1}{Proved}{The displayed fixed-depth outputs tend to the rational pair}
\label{prop:ghost-a65-limit}
The abelian pairs prescribed by the displayed congruence formulas
in SECTIONS A.65 converge to \((3,-2)\) in the power-completion
topology as the depth tends to infinity.
\end{result}
\begin{proof}
The fixed-depth construction in SECTIONS A.65 has
\(D=(R!)^{h_R+1}\mid M_1\mid C\) and
\(A/3=A_0^C\), \(B/(-2)=B_0^C\). Its displayed abelian outputs
therefore converge to \((3,-2)\) as \(R\to\infty\): for a
fixed \(N\), all \(R\ge N\) have \(N\mid R!\mid C\).
This is an exact check of the algorithm's congruences. The
existence or impossibility of a different compatible infinite
family is not determined by this calculation.
\end{proof}

\subsection{An explicit first failure of local matching at class two}
\begin{result}{C.2.2}{Proved}{A factorial ordinary pair first excluded by local matching at class two}
\label{thm:ghost-factorial-selmer32}
There is a completely specified variant of the prime recursion of
SECTIONS A.44, of the same form
\[
 A=3\prod_{j\ge0}q_j^{(j+4)!},\qquad
 B=-2\prod_{j\ge0}r_j^{(j+4)!},
\]
whose first failure of everywhere locally point-or-cusp-matching
nilpotent lifting is class two. The obstructing finite geometric
quotient has order \(32\). This concerns local matching; it
assigns no failure to unrestricted global class-two lifting.
\end{result}

\begin{proof}
The complete radical-field, CRT, Hensel and local-limit proof of
A.44 is reconstructed in the accompanying audit. In its first
stage, choose the old unit solution at \(p=1129\) with
\(X=2\), \(Y\equiv56\pmod p\). The finite identities
\[
 2^{24}\equiv276,\qquad56^{24}\equiv978,\qquad
 3\cdot276-2\cdot978\equiv1\pmod{1129}
\]
give a simple-root Hensel lift. This prime is among the first
old primes and does not divide the next exponent \(120\), so
this choice meets every condition of the recursion. At subsequent
stages choose \(X=1\) and the root \(Y\) near one. The
maintained sum-error inequality supplies precisely the root
precision required in A.44, including the stages at which
\(1129\) divides the factorial. All other choices are the least
eligible choices in that construction. The resulting ordinary
local point satisfies
\[
       t_{1129}\equiv828\pmod{1129},
       \qquad (A,B)\bmod8=([3],[-2]).
\]
The first equality follows at the first update, and every later
multiplier at this prime is one modulo \(1129\). The infinite
support and all other local unit equations are preserved by the
same proved recursion. No numerical execution of the immense
Chebotarev prime search is claimed.

Consider the actual cover and its proper genus-two model
\[
 W:\ z^8=-t(1-t)^3/24,
 \qquad C:\ y^2=x^5-54x,
\]
with mutually inverse function-field formulas
\[
 t=\frac{x^4}{x^4-54},\quad z=\frac{3x}{y},\qquad
 x=-\frac{6z^2}{1-t},\quad y=-\frac{18z}{1-t}.
\]
Over zero and one the cover has one point of ramification index
eight; over infinity it has four points of index two. Thus
\(2g-2=-16+7+7+4=2\), and its removed boundary has six points.
Let \(H\subset\Pi\) be its geometric index-eight subgroup and
let \(K\) be the kernel of the actual compactification and
Abel--Jacobi map
\[
 H\longrightarrow\pi_1(C_{\overline{\mathbb Q}})^{\rm ab}
       \longrightarrow J[2],\qquad J=\operatorname{Jac}(C).
\]
The origin is the rational point at infinity. Compactification and
the mod-two kernel are stable under the deck and arithmetic
actions, so \(K\) is normal in \(E\).

For an exact group computation, let the geometric peripheral
generators be \(x_0,y_0\), and put \(v=x_0\),
\(u=y_0x_0^{-3}\). Schreier generators for \(H\) are
\(u_j=v^juv^{-j}\), \(0\le j<8\), and \(z_0=v^8\).
The zero-boundary loop kills \(z_0\). The one-boundary loop
\(y_0^8\) has homology \(\sum_j u_j+3z_0\). Since
\(x_0y_0=u_1v^4\), its square is \(u_1u_5z_0\).
Its four conjugates give exactly the four infinity-boundary
relations \(u_{j+4}=u_j\). The one-boundary relation is then
redundant modulo two. Hence
\[
 \Pi/K\simeq\bigl(\mathbb F_2[\delta]/\delta^4\bigr)
      \rtimes C_8,\qquad\delta=v-1.
\]
The image of \(v\) has order eight since its eighth power is
the killed cusp loop. Quotienting by \(\delta^2J[2]\) gives
\[
       \overline Q=(\mathbb F_2)^2\rtimes C_8,
       \qquad |\overline Q|=32,
\]
where the generator exchanges the two coordinates. Its
commutator subgroup is the central nonzero diagonal \(C_2\).
Therefore its class is exactly two, and we have the actual
arithmetic quotient map \(E_2\to\overline E\).

Write \(r^2=6\), \(L=\mathbb Q(r)\),
\(\theta^2=3r\), \(K_4=\mathbb Q(\theta)\). The four roots
\(i^j\theta\) give the branch differences
\(T_j=[(i^j\theta,0)-\infty]\); put
\(T=[(0,0)-\infty]\). The basis \(e_j=T+T_j\) satisfies
\(e_2(e_j,T_k)=(-1)^{\delta_{jk}}\) and is cycled by the deck
action. Pairing opposite roots gives a Galois-equivariant
isomorphism
\[
 \overline M=J[2]/\delta^2J[2]
    \simeq\operatorname{Res}_{L/\mathbb Q}\mu_2.
\]
The induced continuous cohomology map is exactly
\[
 K_4^\times/K_4^{\times2}
     \longrightarrow L^\times/L^{\times2},
       \qquad[a]\longmapsto[N_{K_4/L}(a)].
\]
The point-divisor formula is
\([P-\infty]\mapsto[x(P)^2-3r]\). These identifications
follow from \(\operatorname{div}(x-\theta_j)=2T_j\) and
the stated dual Weil basis; the usual odd-degree Kummer map is
recorded in \cite[\S4]{GhostGross2012}. The rational fifth branch
coordinate is reconstructed by the norm, so there is no additional
scalar quotient. In particular
\[
 D=[(3,9)-\infty]\longmapsto[9-3r],\qquad
 T\longmapsto[-3r].
\]

Define \(\mathcal S\subset L^\times/L^{\times2}\) by
membership at every place in the projected local Kummer image of
\(J(\mathbb Q_w)/2J(\mathbb Q_w)\). We compute the entire
group. Good reduction outside two and three and properness show
that these local classes are unramified away from those primes.
The ring of integers of \(L\) is \(\mathbb Z[r]\), with
discriminant \(24\). The Minkowski bound
\(\sqrt{24}/2<3\), together with the principal prime
\(\pi_2=2+r\) of norm \(-2\), proves that its class group
is zero \cite[Theorem 4.3]{GhostMilneANT}. Set
\(\pi_3=3+r\), \(\epsilon=5+2r\). The unit group is
\(\langle-1,\epsilon\rangle\): norm minus one is impossible
modulo three, and a positive norm-one unit between one and
\(\epsilon\) would have coefficient one of \(r\), forcing
\(a^2=7\). Reduction by powers of \(\epsilon\) completes
the proof. We have
\(\pi_2^2=2\epsilon\), \(\pi_3^2=3\epsilon\).
Thus all possible supported global squareclasses are uniquely
\[
         (-1)^a\epsilon^b\pi_2^c\pi_3^d,
               \qquad a,b,c,d\in\mathbb F_2.
\]
At the real embedding \(r=-\sqrt6\), the norm from \(K_4\)
is positive, forcing \(a=c\). At three, \(X^4-54\) is
irreducible: the valuation of a factor's constant term of degree
\(j\) would be \(3j/4\), so \(4\mid j\). Thus
\(J(\mathbb Q_3)[2]\) has order two. The formal pro-three
subgroup has invertible multiplication by two; its finite
quotient consequently gives
\(|J(\mathbb Q_3)/2J(\mathbb Q_3)|=2\). The class
\([-3r]\) has odd valuation and is nonzero. The unit
\(-\pi_2\) is a square at three, \(\epsilon\) has nonsquare
residue, and \(\pi_3\) has odd valuation. This forces \(b=0\).
Finally
\[
       9-3r=9/\pi_3,\qquad
       \frac{-3r}{9-3r}=-\pi_2
\]
shows that the actual rational classes \(D,T\) realize every
survivor and hence every local condition, including the dyadic
one. Therefore
\[
              \mathcal S=\langle[-\pi_2],[\pi_3]\rangle
                    \simeq(\mathbb Z/2)^2.
\]

At \(1129\), the two values of \(r\) are \(838,291\), and
the two generators are squares at both embeddings:
\[
\begin{array}{c|cc}
 &r=838&r=291\\\hline
 -\pi_2&289=17^2&836=426^2\\
 \pi_3&841=29^2&294=221^2
\end{array}\pmod{1129}.
\]
The radicand at \(t=828\) is \(1115\), with exactly the eight
roots \(35,\allowbreak44,\allowbreak235,\allowbreak511,\allowbreak618,\allowbreak894,\allowbreak1085,\allowbreak1094\). Their four
\(x\)-values are \(28,941,1101,188\), each appearing twice.
For all eight sheets the two norm residues are
\((528,1040)\) or \((89,601)\), and all four entries are
nonsquares. Their local Kummer vector is therefore \((1,1)\),
whereas the localized global group \(\mathcal S\) has only
\((0,0)\). All roots and norm residues are units and the root
derivatives are units, so the same statements hold throughout
\(828+1129\mathbb Z_{1129}\).

A locally point-or-cusp-matching \(\overline E\)-section in
combined eighth-power class \([-24]\) trivializes the pullback
of \(W\) and selects a rational sheet. Compactification,
Abel--Jacobi and the quotient \(\overline M\) give a global
class in \(\mathcal S\). Each local geometric conjugator
transports that same sheet. Ordinary sheets yield rational cover
points; cuspidal sheets compactify to rational boundary points.
Every possible transported sheet at \(1129\) has just been
excluded, a contradiction. Our pair has combined class
\(AB^3=[-24]\) and the specified local centre, so its locally
matching class-two lift is excluded. Its abelian section is
locally matching by construction, including the real sign pair
of the point three. This proves exact first failure at class two.

The program \texttt{check\_ghost\_selmer32.py} verifies all
six boundary-loop words, all elements and lower-central orders of
the groups of orders \(128\) and \(32\), their quotient map,
the sixteen supported squareclass conditions, and every one of
the eight local sheets. The complete output is retained.
\end{proof}

\subsection{The ordinary-point scope of SECTIONS A.196}
\begin{result}{C.2.3}{Proved}{The displayed A.196 pair fails the ordinary local test}
\label{prop:ghost-a196-scope}
In \(k=\mathbb Q(\alpha)\), \(\alpha^3=3\), set
\(A=649^8\alpha\), \(U=\alpha^2-2\), \(B=U^3\).
Its image modulo ninth powers at the place above three belongs
to no ordinary or cuspidal point image.
\end{result}
\begin{proof}
Normalize \(v(\alpha)=1\). Then
\[
 U-1=\alpha^2(1-\alpha),\qquad B-1=-6\alpha^2U,
 \qquad v(A)=1, v(U-1)=2, v(B-1)=5.
\]
A unit \(Y\equiv1\pmod\alpha\) satisfies
\(v(Y^3-1)\ge3\) and \(v(Y^9-1)\ge6\), by expansion.
An ordinary ninth-power point would satisfy
\(AX_0^9+BX_1^9=1\). Its first coordinate has valuation
\(1\bmod9\). A negative valuation would force equality with
the second valuation, which is \(0\bmod9\), a contradiction.
A positive valuation forces \(X_1\) to be a unit of residue one
in \(\mathbb F_3\), and then
\(v(1-BX_1^9)=5\), again incompatible with \(1\bmod9\).
Also \(B\) is not a ninth power: the local cubic field contains
no \(\mu_3\), so a ninth root of \(U^3\) would force a cube
root of \(U\), contradicting the valuation two of \(U-1\).
The three cusp conditions respectively require trivial \(B\),
trivial \(A\), or equal two valuations, and all fail.

The independent program \texttt{check\_a196\_local\_scope.py}
enumerates all \(729\) elements of
\((\mathbb Z/9)[\alpha]/(\alpha^3-3)
=\mathcal O_{k_3}/(\alpha^6)\), including all \(243\) units
congruent to one modulo \(\alpha\), and verifies the integral
ninth-power obstruction. Negative valuations were treated above.
No global nonvanishing assertion is inferred from this calculation.
\end{proof}
